\scanpage{133}
Notons que $\lambda'=\lambda^{1/i}$ sera aussi dans $\mathbb L$ (car $\lambda\bar\lambda=Q\,\unsure{\varepsilon|\in|i}$, et
de $=$ pour ses conjugués\ldots), donc on voit aussi qu'il
existe un $\lambda'$ \deleted{\unsure{entier}} $\in\mathbb L$, avec $\deleted{\unsure{\lambda}}\lambda'\bar\lambda'=|\lambda'|^2=Q$, tel que
$\mathbb Q(\lambda')$ soit \deleted{\unsure{pu|p}} \inserted{(purement)} imaginaire et \uline{contienne} $K$; de plus, si
$Q$ est entier, on peut prendre $\lambda'$ entier.
% paper watermark in outline at centre; paper clip at upper left
% archive sheet number: 132
